\paragraph{确定性面积律：\texorpdfstring{\(\det(L_k+tI)\)}{det(L_k+tI)} 的 \(q\)-形式、Szeg\H{o} 分解与 resolvent/Green 闭式}
本段把推论 \ref{cor:pom-Lk-shifted-det-free-energy} 的行列式闭式进一步改写为 \(q\)-形式，并给出与符号积分完全对齐的 bulk--surface 分解，
以及 resolvent 迹与 Green 核矩阵元的双曲函数闭式。

\begin{corollary}[确定性“面积律”分解：\(q\)-形式与指数级有限尺寸余项]\label{cor:pom-Lk-area-law-qform}
\leanverified{paper\_pom\_lk\_area\_law\_qform}
取 \(t>0\)，并定义
\[
\eta(t):=\mathrm{arsinh}\sqrt{\frac{t}{4}}\qquad(\Re\eta\ge 0),
\qquad
q(t):=e^{-2\eta(t)}\in(0,1).
\]
则存在严格恒等式
\[
\det(L_k+tI)=\frac{\cosh\!\bigl((2k+1)\eta(t)\bigr)}{\cosh(\eta(t))}
=q(t)^{-k}\,\frac{1+q(t)^{2k+1}}{1+q(t)}.
\]
从而
\[
\log\det(L_k+tI)=2k\,\eta(t)-\log(1+q(t))+\log\bigl(1+q(t)^{2k+1}\bigr),
\]
其中最后一项为完全显式的有限尺寸余项，并满足
\[
0<\log\bigl(1+q(t)^{2k+1}\bigr)<q(t)^{2k+1}.
\]
\end{corollary}
\begin{proof}
由推论 \ref{cor:pom-Lk-shifted-det-free-energy}，令 \(x=\sqrt{1+t/4}=\cosh(\eta(t))\) 则
\[
\det(L_k+tI)=\frac{T_{2k+1}(x)}{x}=\frac{\cosh((2k+1)\eta(t))}{\cosh(\eta(t))}.
\]
再令 \(q=e^{-2\eta}\)，把 \(\cosh((2k+1)\eta)/\cosh(\eta)\) 写成指数形式即可得到
\(
q^{-k}(1+q^{2k+1})/(1+q)
\)；
取对数即得所示分解。
最后，由 \(q\in(0,1)\) 知 \(\log(1+q^{2k+1})\in(0,q^{2k+1})\)。证毕。
\end{proof}

\begin{proposition}[强 Szeg\H{o} 型分解的可积闭式：符号 \(a_t(\theta)=1+\frac{t}{4\sin^2(\theta/2)}\)]\label{prop:pom-Kk-strong-szego-bulk-surface}
\leanverified{paper\_pom\_kk\_strong\_szego\_bulk\_surface}
由 \(\det(K_k)=1\)（引理 \ref{lem:pom-Kk-gram-det}）知 \(\det(I+tK_k)=\det(L_k+tI)\)。
对 \(t>0\) 定义符号函数
\[
a_t(\theta):=1+\frac{t}{4\sin^2(\theta/2)}\qquad(\theta\in(0,\pi)).
\]
则存在严格分解
\[
\log\det(I+tK_k)
=k\cdot \frac{1}{\pi}\int_{0}^{\pi}\log a_t(\theta)\,\dd\theta
\;+\;\log\frac{1+q(t)^{2k+1}}{1+q(t)},
\]
并且体相积分项可显式化为
\[
\frac{1}{\pi}\int_{0}^{\pi}\log\!\Bigl(1+\frac{t}{4\sin^2(\theta/2)}\Bigr)\,\dd\theta
=2\,\eta(t).
\]
因此 \(\log\det(I+tK_k)\) 的 bulk 密度恰为 \(2\eta(t)\)，而符号中分母的系数 \(4\) 由该闭式同构与 Chebyshev--双曲参数 \(\eta\) 的对齐结构性强制固定。
\end{proposition}
\begin{proof}
由推论 \ref{cor:pom-Lk-area-law-qform}，
\(
\log\det(I+tK_k)=2k\,\eta(t)+\log\frac{1+q^{2k+1}}{1+q}
\)。
故只需证明
\(
\frac{1}{\pi}\int_0^\pi \log a_t(\theta)\,\dd\theta=2\eta(t)
\)。
设
\(
J(t):=\frac{1}{\pi}\int_0^\pi \log\bigl(1+\frac{t}{4\sin^2(\theta/2)}\bigr)\,\dd\theta
\)。
对 \(t>0\) 可在积分号下求导得
\[
J'(t)=\frac{1}{\pi}\int_0^\pi \frac{1}{t+4\sin^2(\theta/2)}\,\dd\theta
=\frac{1}{\pi}\int_0^\pi \frac{1}{t+2(1-\cos\theta)}\,\dd\theta
=\frac{1}{\sqrt{t(4+t)}},
\]
其中最后一步为标准积分恒等式（亦可由复积分或对 \(t\) 的 Green 函数解析解得到）。
另一方面，\(\eta(t)=\mathrm{arsinh}\sqrt{t/4}\) 满足 \(\eta'(t)=1/(2\sqrt{t(4+t)})\)，故 \((2\eta)'=J'\)。
又 \(J(0)=0=2\eta(0)\)，从而 \(J(t)=2\eta(t)\)。证毕。
\end{proof}

\begin{corollary}[解析可审计的 resolvent 迹：\texorpdfstring{\(\tr(L_k+tI)^{-1}\)}{tr((L_k+tI)^{-1})} 的双曲闭式]\label{cor:pom-Lk-resolvent-trace-hyperbolic}
取 \(t>0\) 并沿用 \(\eta(t)\) 记号。
则 resolvent 迹满足严格恒等式
\[
\tr(L_k+tI)^{-1}
=\frac{\partial}{\partial t}\log\det(L_k+tI)
=\frac{2k+1}{2\sqrt{t(4+t)}}\,\tanh\!\bigl((2k+1)\eta(t)\bigr)-\frac{1}{2(4+t)}.
\]
特别地，迹密度极限存在且为
\[
\lim_{k\to\infty}\frac{1}{2k+1}\tr(L_k+tI)^{-1}
=\frac{1}{2\sqrt{t(4+t)}}.
\]
\end{corollary}
\begin{proof}
由推论 \ref{cor:pom-Lk-shifted-det-free-energy}，
\(
\log\det(L_k+tI)=\log\cosh((2k+1)\eta)-\log\cosh\eta
\)（其中 \(\eta=\eta(t)\)）。
对 \(t\) 求导并用 \(\eta'(t)=1/(2\sqrt{t(4+t)})\) 得
\[
\frac{\partial}{\partial t}\log\det(L_k+tI)
=\eta'(t)\Bigl((2k+1)\tanh((2k+1)\eta)-\tanh\eta\Bigr)
=\frac{2k+1}{2\sqrt{t(4+t)}}\tanh((2k+1)\eta)-\frac{1}{2(4+t)},
\]
其中最后一步用到 \(\tanh(\eta)=\sqrt{t/(4+t)}\)。
而 \(\tr(L_k+tI)^{-1}=\partial_t\log\det(L_k+tI)\) 为一般恒等式。
取 \(k\to\infty\) 且用 \(\eta>0\Rightarrow\tanh((2k+1)\eta)\to 1\) 即得极限。证毕。
\end{proof}

\begin{proposition}[Green 核矩阵元的完全闭式：\texorpdfstring{\(((L_k+tI)^{-1}_{ij})\)}{((L_k+tI)^{-1}_{ij})} 的双曲乘积公式]\label{prop:pom-Lk-green-kernel-entries}
\leanverified{paper\_pom\_Lk\_green\_kernel\_entries}
取 \(t>0\) 并记 \(\eta=\eta(t)\)。
则对任意 \(1\le i\le j\le k\)，逆矩阵矩阵元满足严格闭式
\[
(L_k+tI)^{-1}_{ij}
=\frac{\sinh(2i\eta)\,\cosh\!\bigl((2(k-j)+1)\eta\bigr)}
{\sinh(2\eta)\,\cosh\!\bigl((2k+1)\eta\bigr)},
\qquad
(L_k+tI)^{-1}_{ji}=(L_k+tI)^{-1}_{ij}.
\]
特别地，
\[
(L_k+tI)^{-1}_{11}=\frac{\cosh((2k-1)\eta)}{\cosh((2k+1)\eta)},
\qquad
(L_k+tI)^{-1}_{1k}=\frac{\cosh(\eta)}{\cosh((2k+1)\eta)}.
\]
\end{proposition}
\begin{proof}
令 \(A:=L_k+tI\)。它是对称三对角矩阵，除末端外对角元恒为 \(2+t\)，末端对角元为 \(1+t\)，非对角元恒为 \(-1\)。
对 \(r\ge 0\) 定义 leading 主子式
\(
D_r:=\det(A_{[1..r,\,1..r]})
\)
并约定 \(D_0:=1\)。则对 \(1\le r\le k-1\) 有递推 \(D_r=(2+t)D_{r-1}-D_{r-2}\)，从而
\[
D_{i-1}=\frac{\sinh(2i\eta)}{\sinh(2\eta)}.
\]
同理，定义 trailing 主子式
\(
E_r:=\det(A_{[k-r+1..k,\,k-r+1..k]})
\)
并约定 \(E_0:=1\)、\(E_1:=1+t\)。递推与端点条件解出
\[
E_{k-j}=\frac{\cosh\!\bigl((2(k-j)+1)\eta\bigr)}{\cosh(\eta)}.
\]
对三对角矩阵的 Cramer--Gantmacher 公式（亦即一维差分 Green 函数的标准表示）给出：当 \(i\le j\) 时
\[
A^{-1}_{ij}=\frac{D_{i-1}\,E_{k-j}}{D_k},
\]
其中 \(D_k=\det(A)=\cosh((2k+1)\eta)/\cosh(\eta)\) 由推论 \ref{cor:pom-Lk-shifted-det-free-energy} 给出。
代入即得所示闭式。证毕。
\end{proof}

\begin{corollary}[指数聚类与相关长度：质量参数 \(t\) 的严格“间隙律”]\label{cor:pom-Lk-exponential-clustering-gaplaw}
\leanverified{paper\_pom\_lk\_exponential\_clustering\_gaplaw}
取 \(t>0\) 并令 \(\eta=\eta(t)>0\)。则对任意 \(1\le i\le j\le k\) 有一致上界
\[
(L_k+tI)^{-1}_{ij}\ \le\ \frac{1}{\sinh(2\eta)}\,e^{-2\eta(j-i)}.
\]
从而可定义相关长度
\[
\xi(t):=\frac{1}{2\eta(t)}=\frac{1}{2\,\mathrm{arsinh}\sqrt{t/4}},
\]
并得到严格指数聚类：在固定 \(t>0\) 下，Green 关联随格距 \(j-i\) 以速率 \(e^{-(j-i)/\xi(t)}\) 衰减。
\end{corollary}
\begin{proof}
由命题 \ref{prop:pom-Lk-green-kernel-entries}，
\[
(L_k+tI)^{-1}_{ij}
=\frac{\sinh(2i\eta)\,\cosh((2(k-j)+1)\eta)}{\sinh(2\eta)\,\cosh((2k+1)\eta)}.
\]
用 \(\sinh u\le \tfrac12 e^{u}\) 与 \(\cosh u\le e^{u}\)，以及 \(\cosh((2k+1)\eta)\ge \tfrac12 e^{(2k+1)\eta}\)，得到
\[
(L_k+tI)^{-1}_{ij}
\le
\frac{(\tfrac12 e^{2i\eta})\cdot (e^{(2(k-j)+1)\eta})}{\sinh(2\eta)\cdot (\tfrac12 e^{(2k+1)\eta})}
=\frac{1}{\sinh(2\eta)}\,e^{-2\eta(j-i)}.
\]
相关长度表述由 \(\xi=1/(2\eta)\) 的重写立得。证毕。
\end{proof}

\begin{corollary}[半无限极限的边界固定点：反射系数 \(q(t)\) 的代数闭式与 Riccati 方程]\label{cor:pom-Lk-boundary-fixed-point-riccati}
\leanverified{paper\_pom\_lk\_boundary\_fixed\_point\_riccati}
取 \(t>0\) 并令 \(q(t)=e^{-2\eta(t)}\)。
则边界 Green 值存在且满足
\[
\lim_{k\to\infty}(L_k+tI)^{-1}_{11}=q(t).
\]
并且 \(q(t)\) 具有闭式
\[
q(t)=\Bigl(\frac{\sqrt{4+t}-\sqrt{t}}{2}\Bigr)^{2}
=1+\frac{t}{2}-\frac{1}{2}\sqrt{t(4+t)},
\]
且是二次方程的唯一 \((0,1)\) 解：
\[
q^2-(2+t)q+1=0.
\]
\end{corollary}
\begin{proof}
由命题 \ref{prop:pom-Lk-green-kernel-entries}，\((L_k+tI)^{-1}_{11}=\cosh((2k-1)\eta)/\cosh((2k+1)\eta)\)，令 \(k\to\infty\) 得极限为 \(e^{-2\eta}=q\)。
又 \(\eta=\mathrm{arsinh}(\sqrt{t}/2)\)，故 \(e^{-\eta}=\sqrt{1+t/4}-\sqrt{t}/2\)，从而
\[
q=e^{-2\eta}=\Bigl(\sqrt{1+\frac{t}{4}}-\frac{\sqrt{t}}{2}\Bigr)^2
=\Bigl(\frac{\sqrt{4+t}-\sqrt{t}}{2}\Bigr)^2,
\]
展开即得第二式。
最后，由 \(2\cosh(2\eta)=e^{2\eta}+e^{-2\eta}=q^{-1}+q=2+t\) 得 \(q^2-(2+t)q+1=0\)，并由 \(q\in(0,1)\) 的分支选取唯一性得结论。证毕。
\end{proof}

\begin{corollary}[表面自由能与边界响应：常数项闭式与可观测导数]\label{cor:pom-Lk-surface-free-energy}
\leanverified{paper\_pom\_lk\_surface\_free\_energy}
取 \(t>0\) 并记 \(\eta=\eta(t)\)、\(q=q(t)\)。
定义 bulk 自由能密度与表面自由能
\[
f(t):=\lim_{k\to\infty}\frac{1}{k}\log\det(L_k+tI)=2\eta(t),
\qquad
g(t):=\lim_{k\to\infty}\Bigl(\log\det(L_k+tI)-k f(t)\Bigr)=-\log(1+q(t)).
\]
则对任意 \(k\ge 1\) 有严格恒等式
\[
\log\det(L_k+tI)=k f(t)+g(t)+\log\bigl(1+q(t)^{2k+1}\bigr),
\]
且表面项的导数可由边界响应给出：
\[
g'(t)=\lim_{k\to\infty}\Bigl(\tr(L_k+tI)^{-1}-k f'(t)\Bigr)
=\frac{q(t)}{1+q(t)}\cdot \frac{1}{\sqrt{t(4+t)}}.
\]
\end{corollary}
\begin{proof}
由推论 \ref{cor:pom-Lk-area-law-qform} 的精确分解
\(
\log\det=2k\eta-\log(1+q)+\log(1+q^{2k+1})
\)，取极限即得 \(f(t)=2\eta\) 与 \(g(t)=-\log(1+q)\) 及所示恒等式。
对导数部分，注意 \(f'(t)=2\eta'(t)=1/\sqrt{t(4+t)}\)，并由推论 \ref{cor:pom-Lk-resolvent-trace-hyperbolic} 得
\[
\tr(L_k+tI)^{-1}-k f'(t)
=\frac{2k+1}{2\sqrt{t(4+t)}}\tanh((2k+1)\eta)-\frac{1}{2(4+t)}-\frac{k}{\sqrt{t(4+t)}}.
\]
令 \(k\to\infty\)（\(\eta>0\Rightarrow\tanh((2k+1)\eta)\to 1\)）得到极限
\(
\frac{1}{2\sqrt{t(4+t)}}-\frac{1}{2(4+t)}
\)。
而由推论 \ref{cor:pom-Lk-boundary-fixed-point-riccati} 的 \(q\)-闭式可化简为
\(
\frac{q}{1+q}\cdot \frac{1}{\sqrt{t(4+t)}}
\)，与 \(g'(t)=-(q'/(1+q))\)（\(q'=-(2\eta')q\)）一致。证毕。
\end{proof}

\begin{corollary}[表面自由能的严格凹性：二阶导数闭式与符号刚性]\label{cor:pom-Lk-surface-free-energy-strict-concavity}
\leanverified{paper\_pom\_lk\_surface\_free\_energy\_strict\_concavity}
沿用推论 \ref{cor:pom-Lk-surface-free-energy} 的记号。
则对任意 \(t>0\) 有二阶导数闭式
\[
\boxed{\;
g''(t)
=-\frac{q(t)}{(1+q(t))^2\,t(4+t)}
-\frac{q(t)(t+2)}{(1+q(t))\,[t(4+t)]^{3/2}}\;<\;0\;}
\]
其中 \(q(t)=1+\frac{t}{2}-\frac12\sqrt{t(4+t)}\in(0,1)\) 为边界固定点（推论 \ref{cor:pom-Lk-boundary-fixed-point-riccati}）。
因此 \(g\) 在整个正轴上严格凹，其凹性由 \(q(t)\) 的代数闭式强制。
\end{corollary}
\begin{proof}
由 \(g(t)=-\log(1+q(t))\) 与推论 \ref{cor:pom-Lk-surface-free-energy} 的导数闭式
\(
g'(t)=\frac{q(t)}{1+q(t)}\cdot\frac{1}{\sqrt{t(4+t)}}
\)
直接求导，并用 \(q(t)=1+\frac{t}{2}-\frac12\sqrt{t(4+t)}\) 化简，即得所示表达式。
由于 \(q(t)\in(0,1)\) 且 \(t(4+t)>0\)，右端两项均严格为负，故 \(g''(t)<0\)。证毕。
\end{proof}

\begin{remark}[有限窗审计：确定性面积律分解、resolvent 迹与 Green 核矩阵元]\label{rem:pom-fence-green-kernel-area-law-audit}
脚本在有限窗上核验：推论 \ref{cor:pom-Lk-area-law-qform} 的 \(q\)-形式分解与指数小余项界、
命题 \ref{prop:pom-Kk-strong-szego-bulk-surface} 的符号积分闭式、
推论 \ref{cor:pom-Lk-resolvent-trace-hyperbolic} 的 resolvent 迹闭式、
命题 \ref{prop:pom-Lk-green-kernel-entries} 的矩阵元闭式与推论 \ref{cor:pom-Lk-exponential-clustering-gaplaw} 的指数聚类上界，
并对边界固定点 \(q(t)\) 与表面自由能导数 \(g'(t)\) 给出数值一致性证书。
审计表如下：
\begin{table}[H]
\centering
\scriptsize
\setlength{\tabcolsep}{6pt}
\caption{确定性面积律/Green-resolvent 闭式的有限窗审计（取 $t=1$）：核验 $q$-形式行列式分解（推论~\ref{cor:pom-Lk-area-law-qform}）、符号积分闭式（命题~\ref{prop:pom-Kk-strong-szego-bulk-surface}）、resolvent 迹闭式（推论~\ref{cor:pom-Lk-resolvent-trace-hyperbolic}）、Green 矩阵元闭式与指数聚类上界（命题~\ref{prop:pom-Lk-green-kernel-entries}、推论~\ref{cor:pom-Lk-exponential-clustering-gaplaw}），以及边界固定点 $q(t)$ 与表面项导数 $g'(t)$ 的一致性（推论~\ref{cor:pom-Lk-boundary-fixed-point-riccati}、\ref{cor:pom-Lk-surface-free-energy}）。}
\label{tab:fence_green_kernel_area_law_resolvent_audit}
\begin{tabular}{r r r r r r r r}
\toprule
$k$ & $|\Delta\log\det|$ & $|\Delta\,\mathrm{bulk}|$ & $|\Delta\,\mathrm{tr}|$ & $\max|\Delta G_{ij}|$ & $\max\rho$ & $|G_{11}-q|$ & $|\Delta g'|$\\
\midrule
10 & 0.00e+00 & 1.20e-04 & 8.88e-16 & 2.78e-16 & 0.590170 & 3.33e-16 & 3.91e-15\\
20 & 3.55e-15 & 1.20e-04 & 0.00e+00 & 3.89e-16 & 0.676275 & 3.33e-16 & 3.91e-15\\
40 & 7.11e-15 & 1.20e-04 & 3.55e-15 & 5.55e-16 & 0.500000 & 3.33e-16 & 3.91e-15\\
200 & 0.00e+00 & 1.20e-04 & 0.00e+00 & -- & 0.500000 & 3.33e-16 & 7.47e-15\\
\bottomrule
\end{tabular}
\end{table}

\end{remark}
